\chapter{Abril}

\section{Fiktionsbescheinigung - 01}
Quarta-feira pegamos carona com Uwe do departamento de Recursos Humanos para o escritório de imigração realizar o documento de permissão da residência na Alemanha pelo período estabelecido no contrato.

\section{Segurança do Trabalho - 07}
Na terça-feira após o feriado participamos da palestra sobre segurança do trabalho apresentada por Menno, onde recebemos as instruções básicas de segurança e evacuação. Após a palestra, recebemos os equipamentos de segurança e seguimos para as salas de escritório onde Luciano nos entregou instruções de login nos computadores. Após me estabelecer na mesa e realizar o login, não haviam atividades posteriores, visto que os grupos e coordenadores seriam discutidos no dia seguinte por Benjamin, então me encaminhei de volta para casa.

\section{Divisão de Grupos - 08}
Nesse dia, Benjamin nos apresentou o departamento e os projetos em andamento. Após ter explicado os detalhes de cada projeto, discutimos quais pessoas contribuiriam melhor para cada tipo de trabalho. 

\begin{itemize}
	\item Friction Spot Welding
	\item Bobbin Tools
	\item Friction Extrusion
\end{itemize}

Durante a conversa sobre a divisão dos grupos, demostrei meu interesse na área de simulação numérica que vi presente em alguns dos projetos, Benjamin achou interessante e me chamou para conversar mais sobre o assunto após a apresentação. Nessa conversa, ele mostrou ou projetos que faziam o uso dessa área e precisavam de ajuda em alguns aspectos. Ele me mostrou artigos desenvolvidos e me recomendou leitura de um deles para ambientação. Após a conversa, definimos que por enquanto ficarei mais em contato com a Camila no projeto de Spot Welding, pois existem alguns trabalhos relacionados que entram na área de simulação e coleta de dados. Por fim, falei com Menno para levar meu computador para o serviço de TI, pois o armazenamento estava cheio, então sexta-feira ele via levar para o serviço.

\section{Projeto de Spot Welding - 09}
Logo pela manhã tive o primeiro contato com Camila, que me apresentou os projetos dos quais ela é responsável. Ficou acordado que por enquanto ajudarei na parte experimental do Projeto de Spot Welding e semana que vem, Benjamin discutirá comigo mais sobre projetos que envolvam propagação de trinca (FCP - Fracture Crack Propagation) e simulação numérica.

Sobre o computador na minha mesa no escritório, Menno passou na sala pela manha para limpar o armazenamento mas precisava do acesso de administrador. Portanto, acabou levando o computador ao setor de TI para limpar e colocar um HD de armazenamento extra. Por enquanto estou usando meu notebook pessoal.

Mais tarde tivemos o treinamento para realizar ensaios de tração com o Fares.

\section{Treinamento de Metalografia - 10}
Nesse dia tivemos o treinamento de metalografia e logo depois Camila pediu para realizar a análise metalográfica de seis amostras de Spot Welding em alumínio. Para isso a estrutura soldada precisa ser cordada, então o treinamento para máquina de corte foi combinado com Menno para segunda-feira.

\section{Treinamento de Corte - 13}
Foi realizado o treinamento de corte com Menno e no fim do dia fomos adicionados no sistema de reservas. Não prosseguimos com o corte das amostras pois ainda não estávamos no sistema de reservas.

\section{Corte e Embutimento das Amostras - 14}
Após reservar a maquina de corte, Mika e eu cortamos as amostras de Spot Welding (seis para cada um de nós). Após o almoço, reservamos o laboratório de metalografia e realizamos o embutimento e desbaste com lixa 80. Amanhã pela manhã, Camila realizará um treinamento de solda conosco, então o lixamento das amostras provavelmente continuará pela tarde.

\section{Treinamento da Máquina de Spot Welding - 15}
Camila realizou um treinamento na máquina de Refill Stir Friction Spot Welding (RSFSW) em que realizamos quatro soldas variando parâmetros para observar o impacto na qualidade. O treinamento decorreu do período da manhã até o início da tarde, então logo depois do almoço Mika e eu prosseguimos com o corte e preparo das soldas recém feitas. Nesse dia, as seis amostras que Camila tinha me entregado no dia 10 foram preparadas até a lixa 800, faltando somente a faze de polimento. As amostras do treinamento de Spot welding foram desbastadas com a lixa 80 somente.

\section{Polimento das Amostras - 16}

Foi feito o polimento em dois estágios das seis amostras, desbaste na lixa 800 das quatro amostras do treinamento e posteriormente o seu polimento, deixando assim, todas as amostras prontas para o ataque químico. Como o ataque pode ser feito somente por um estudante de Doutorado ou superior, combinaremos um horário com Camila semana que vem. Enquanto isso, marcamos treinamento no microscópio para segunda feira, assim, poderemos concluir a análise das amostras após o ataque. 

\section{Detalhes do projeto - 17}
Logo pela manhã, Camila realizou o ataque químico de uma das amostras somente para demostração da microscopia e deu instruções de preparo para as análises de algumas amostras específicas. Após o almoço, Mika e eu fomos para a sala da Camila, onde ela nos deu mais detalhes sobre os projetos dos quais estamos ajudando com o trabalho. Um dos detalhes envolvia ajuda com o tratamento de alguns dados de ensaio, então comecei o processo em Python. 

Nesse dia, Camila iria realizar testes de flexão e pediu nossa presença, mas a pessoa que ela dependia para fazer o ensaio não estava disponível, então voltei ao escritório e prossegui com a leitura de artigos que Camila e Benjamin tinham me encaminhado.

\section{Microscópio - 20}
Na manhã desse dia, participamos do treinamento de microscópio e logo depois ja realizamos fotos das amostras polidas. Camila pediu para tirar foto das dez amostras de cada um antes e depois do ataque químico. Após as fotos tiradas, segui o trabalho de tratamento de dados em Python.

\section{Ataque químico - 21}
Camila conseguiu marcar o laboratório para realizar os ataques químicos pela tarde, mas logo pela manhã ja estávamos acompanhando a soldagem das amostras do colega Ali para praticar a montagem, uso e desmontagem da máquina. Antes de realizar o ataque, foi necessário fazer um processo adicional de preparação das amostras, foi preciso expor a face contraria à polida para que houvesse contato elétrico no ataque químico. Então, Camila fez alguns furos nas amostras e lixamos outras em que o metal estava quase exposto. Por fim, voltamos para a soldagem das amostras do Ali e marcamos treinamento da furadeira de bancada para amanhã. Nesse dia também prossegui o tratamento de dados.



